
\section{Kiến trúc phần mềm nhúng dựa trên hệ điều hành FreeRTOS}

Phần mềm nhúng của hệ thống được phát triển trên FreeRTOS để đáp ứng yêu cầu xử lý thời gian thực \cite{ref_2_freertos}. Trong các hệ thống nhúng đơn giản, cấu trúc vòng lặp tuần tự có thể đủ cho việc đọc cảm biến và điều khiển cơ bản. Tuy nhiên, hệ thống AIoT của đồ án phải thực hiện đồng thời nhiều tác vụ có đặc tính thời gian khác nhau, bao gồm thu thập dữ liệu âm thanh, lấy mẫu rung động, xử lý tín hiệu số, suy luận TensorFlow Lite Micro, duy trì kết nối Wi-Fi và truyền bản tin MQTT. Nếu toàn bộ luồng xử lý được đặt trong một vòng lặp duy nhất, độ trễ mạng hoặc tác vụ tính toán nặng có thể làm mất mẫu cảm biến và ảnh hưởng trực tiếp đến sai số tái tạo.

FreeRTOS cho phép chia hệ thống thành các tác vụ độc lập, mỗi tác vụ có mức ưu tiên và chu kỳ thực thi riêng. Thay vì sử dụng mô hình lập trình tuần tự, đồ án triển khai kiến trúc \textit{song song thực sự trên lõi kép}, khai thác hai lõi vật lý của Seeed Studio XIAO ESP32-S3 để phân tách độc lập luồng thu thập dữ liệu và luồng suy luận AI. Các cơ chế nhóm sự kiện, hàng đợi và ghim lõi được sử dụng nhằm đồng bộ dữ liệu, giảm tranh chấp tài nguyên và bảo vệ luồng suy luận khỏi độ trễ bất định của truyền thông mạng.

\subsection{Thiết kế luồng xử lý tín hiệu âm thanh trên lõi số 0}

Tác vụ xử lý âm thanh được gán vào lõi số 0 với mức ưu tiên 10, đảm bảo dòng dữ liệu liên tục và tránh mất mẫu. Đường ống xử lý bao gồm:

\begin{itemize}
    \item \textbf{Thu thập dữ liệu I2S-DMA:} Cơ chế DMA được sử dụng để đọc liên tục tín hiệu âm thanh từ microphone INMP441. Tần số lấy mẫu được thiết lập ở mức 8 kHz. Việc giới hạn tần số lấy mẫu này giúp tiết kiệm dung lượng bộ nhớ đệm và giảm tải khối lượng tính toán FFT, đồng thời vẫn giữ lại được các dải tần số âm học đặc trưng của động cơ (0-4 kHz).
    \item \textbf{Tiền xử lý DSP:} Tác vụ áp dụng hàm cửa sổ Hann và phép biến đổi Fourier nhanh 512 điểm thông qua thư viện \texttt{esp-dsp} \cite{ref_2_espdsp}.
    \item \textbf{Trích xuất đặc trưng Mel:} Năng lượng trên 13 dải lọc Mel được tính toán theo thang logarit để tạo thành vector đặc trưng âm thanh.
\end{itemize}

Sau khi hoàn tất, tác vụ sẽ đẩy dữ liệu vào \texttt{audio\_queue} và bật cờ hiệu \texttt{AUDIO\_READY\_BIT} trong nhóm sự kiện.

\subsection{Thiết kế luồng thu thập rung động và suy luận AI trên lõi số 1}

Lõi số 1 quản lý hai tác vụ với mức ưu tiên khác nhau, đảm nhiệm thu thập rung động và ra quyết định cảnh báo.

\textbf{Tác vụ thu thập rung động (\texttt{vib\_task}).} Tác vụ này được cấp mức ưu tiên cao nhất trong ứng dụng người dùng (mức ưu tiên 15) nhằm đảm bảo độ chính xác cho chu kỳ lấy mẫu. Thông qua hàm định thời \texttt{vTaskDelayUntil}, hệ thống duy trì chu kỳ 1 ms (1000 Hz) cho mỗi lần đọc cảm biến ADXL345 qua bus I2C. Dữ liệu sau đó được xử lý để lấy giá trị RMS và phương sai trục Z, đồng thời bật cờ hiệu \texttt{VIB\_READY\_BIT}.

\textbf{Tác vụ suy luận AI (\texttt{ai\_task}).} Được cấu hình ở mức ưu tiên 5, tác vụ này đảm nhiệm quá trình phân tích và ra quyết định với luồng thực thi như sau:
\begin{enumerate}
    \item \textbf{Đồng bộ hóa:} Đợi nhóm sự kiện (\texttt{AUDIO\_READY\_BIT} \& \texttt{VIB\_READY\_BIT}).
    \item \textbf{Chuẩn hóa hỗn hợp:} Áp dụng \textit{MinMax Scaling} cho 13 đặc trưng âm thanh và \textit{RobustScaler} cho 6 đặc trưng rung động để hạn chế ảnh hưởng của nhiễu ngoại lai trước khi suy luận.
    \item \textbf{Chuyển đổi bộ thông dịch:} Hệ thống đã cấp phát sẵn 3 vùng đệm tensor trên PSRAM. Dựa vào chỉ số phương sai rung động, hệ thống sẽ trỏ đến bộ thông dịch của mô hình tương ứng (GENTLE, STRONG, hoặc SPIN).
    \item \textbf{Suy luận:} Gọi hàm \texttt{Invoke()} của TensorFlow Lite Micro để tính toán sai số tái tạo (MAE).
    \item \textbf{Ra quyết định:} Đối chiếu giá trị MAE với các ngưỡng tĩnh đã được định nghĩa trong \texttt{model\_data.h} để xác định trạng thái bất thường.
\end{enumerate}

\subsection{Cơ chế đồng bộ hóa đa tác vụ}

Để quản lý luồng dữ liệu, hệ thống sử dụng ba cơ chế đồng bộ hóa của FreeRTOS:

\begin{itemize}
    \item \textbf{Nhóm sự kiện:} Đường ống xử lý AI chỉ được kích hoạt khi cả hai luồng dữ liệu (âm thanh và rung động) đều đã thu thập đủ mẫu, giúp loại bỏ tình trạng lệch pha thời gian giữa các cảm biến.
    \item \textbf{Hàng đợi:} Sử dụng \texttt{audio\_queue} để truyền vector đặc trưng và \texttt{mqtt\_queue} để truyền thông báo cảnh báo. Cơ chế hàng đợi giúp các tác vụ giao tiếp an toàn mà không chặn lẫn nhau.
    \item \textbf{Cô lập tác vụ truyền dẫn:} Quá trình gửi dữ liệu MQTTS được tách biệt khỏi luồng suy luận AI và đưa vào hàm \texttt{loop()} nội tại (mức ưu tiên 1). Hệ thống sử dụng phương thức \texttt{mqtt\_smart\_publish()} để gửi bản tin nhịp định kỳ (30 giây) hoặc gửi ngay lập tức khi có sự kiện cảnh báo. Cách thiết kế này tách biệt độ trễ mạng khỏi vòng lặp xử lý AI tại vùng biên, tránh gián đoạn quá trình suy luận.
\end{itemize}

% ==============================================================================
% DANH MỤC TÀI LIỆU THAM KHẢO TRÍCH DẪN (CHI TIẾT CHO MỤC 2.3)
% ==============================================================================
% \bibitem{ref_2_freertos} Real-Time Engineers Ltd., "FreeRTOS Reference Manual," v10.4.0, 2021. [Online]. Available: https://www.freertos.org
% \bibitem{ref_2_espdsp} Espressif Systems, "ESP-DSP Library Programming Guide," 2023. [Online]. Available: https://docs.espressif.com/projects/esp-dsp/
